\begin{tomtat}
\begin{itemize}
  \item CBM ép luồng tính toán qua một tầng mà mỗi đơn vị được huấn luyện để
nhận đúng một khái niệm do con người gán nhãn, rồi dự đoán nhãn chỉ từ tầng ấy.
So với mười bốn chương trước, nó bỏ được hai thứ cùng lúc: không phải hỏi mô
hình ở những đầu vào đã bị sửa, và không phải đi tìm ground truth, vì nhãn khái
niệm nằm sẵn trong bộ dữ liệu. Các đơn vị của tầng thắt là thành phần theo nghĩa
của mục~\ref{sec:ch07-ba-doi-tuong}, chỉ khác ở chỗ chúng được dựng theo tên có
sẵn thay vì được đặt tên sau.
  \item \tn{Rò rỉ}{leakage} là việc mô hình cất thêm thông tin vào chính các giá
trị khái niệm, và đầu phân loại đọc phần thông tin thêm ấy trong khi người giám
sát thì không. Nó đi xuyên qua tầng thắt chứ không đi vòng, nên bảng khái niệm
vẫn đầy đủ về hình thức, và can thiệp đúng vào một khái niệm có thể làm độ chính
xác nhiệm vụ tụt đi. Đọc theo bảng~\ref{tab:ch01-bon-quan-he}, rò rỉ là họ hàng
của định nghĩa~\ref{def:ch01-do-trung-thuc} và không phải nó: cái được trưng ra
là một tầng của chính mô hình chứ không phải một lời giải thích bên ngoài, phần
chênh được định nghĩa so với một tập khái niệm đã gán nhãn nên đổi tập thì đổi
mốc, và không bài nào trong ba bài gọi nó là độ trung thực.
  \item Định nghĩa~\ref{def:ch15-hai-diem-so} đo rò rỉ bằng phần chênh giữa
lượng thông tin chuẩn hóa mà giá trị học được mang và lượng thông tin mà khái
niệm thật mang, về nhãn cho CTL và về khái niệm khác cho ICL. Cần một định nghĩa
theo lối thông tin vì các điểm đánh giá thông thường không thấy rò rỉ: một cặp
mô hình có sáu điểm đánh giá lệch nhau không quá $0.005$ mà
\tn{điểm số can thiệp}{intervention score} lệch $0.301$.
  \item Bảng~\ref{tab:ch15-tuong-quan-diem-so} là chỗ Phần~IV chờ suốt bảy
chương: một dụng cụ đo được đối chiếu với một tiêu chí độc lập. CTL vượt OIS ở
cả tám bộ dữ liệu và có ý nghĩa thống kê ở cả tám; ICL thì thua OIS ở ba bộ, nên
câu của bài rằng cả hai vượt OIS trên mọi bộ dữ liệu là câu nói quá bảng của
chính nó. Ba chỗ phép kiểm định ấy không với tới: tiêu chí dùng làm mốc đặc hiệu
mà không nhạy, phép ước lượng thông tin không so được hai biểu diễn khác số
chiều, và thứ được kiểm định là điểm số rò rỉ chứ không phải chỉ số độ trung
thực. Điều làm nó chạy được là ground truth có sẵn trong bộ dữ liệu, tức đường
thứ tư sau ba đường mà mục~\ref{sec:ch14-ground-truth-ke-tan-cong} đã đếm.
  \item Bốn nguyên nhân gây rò rỉ chia hai hạng. Hai nguyên nhân đầu là siêu
tham số và chọn được: giám sát khái niệm quá nhẹ, và cách mã hóa giàu hơn mức
cần. Hai nguyên nhân sau thì không: tập khái niệm không đủ dự báo nhiệm vụ, và
đầu phân loại quá nghèo để học quan hệ giữa khái niệm và nhãn. Cái giá của lối
thoát nằm ở hạng thứ hai, và chính người đề xuất khung viết ra: khi không kiếm
thêm được nhãn khái niệm thì phải kết luận nhiệm vụ ấy không dùng lối tiếp cận
theo khái niệm được. Ở CEM thì
\tn{rò rỉ khái niệm-nhiệm vụ}{concepts-task leakage} là thứ kiến trúc chừa sẵn
chỗ cho, không phải tai nạn của phép tối ưu, nên hiệu năng của nó đến từ phần
thông tin thêm chứ không từ phần khái niệm; CEM còn thêm một dạng thứ ba,
\tn{rò rỉ căn chỉnh}{alignment leakage}, chỉ xuất hiện ở mô hình được huấn luyện
kèm can thiệp, và đó là lý do hiệu năng khi can thiệp thôi không đo được rò rỉ ở
kiến trúc ấy.
  \item Bài thứ ba đo tách nhau hai tính chất của từng nơ-ron trong một SAE: có
tên đọc được, và ép được đầu ra đi theo cái tên ấy. Tính chất thứ hai là mức
quan trọng theo nghĩa mục~\ref{sec:ch09-doi-tuong-khac-nhau} đã cố định, không
phải độ trung thực. Chỉ 18.84 phần trăm nơ-ron đạt cả hai, và một từ điển
65\,536 đơn vị chỉ nhận được 28.0 phần trăm số khái niệm trong một danh sách hai
mươi nghìn từ thông dụng. Mức tăng \tn{khả năng điều khiển}{steerability} mà bài
báo cáo cho cả ba mô hình là 14.5 phần trăm, trung bình của ba con số; một trong
ba con số ấy là tổng hai chỉ số chứ không phải trung bình của chúng, và thay nó
bằng trung bình thì mức tăng là 12.2 phần trăm. Mọi thứ hạng giữ nguyên. Cả ba
bài đều không dùng chữ \enquote{faithful} một lần nào trong phần thân.
\end{itemize}
\end{tomtat}

\cauhoi
\begin{cauhoids}
  \item Mục~\ref{sec:ch15-ro-ri} lấy ngưỡng tại 1 trên hàm chạy suốt sách và
được nhiệm vụ \enquote{hoặc}. Hãy lấy ngưỡng tại 2 thay vào đó, viết ra nhãn
$y$ ở bốn tổ hợp, rồi tính $H(y)$, $I(c_1,y)$ và $I(c_2,y)$ dưới cùng giả thiết
độc lập. Từ đó cho biết trần của $\mathrm{CTL}_1$ và trần của $\mathrm{CTL}_2$
trên nhiệm vụ mới, và giải thích vì sao hai trần ấy khác nhau trong khi ở nhiệm
vụ \enquote{hoặc} chúng bằng nhau.
  \item Mục~\ref{sec:ch15-dung-cu-duoc-kiem} gọi lập luận thứ hai của bài báo là
phép kiểm nhất quán chứ không phải phép kiểm định. Hãy dựng một điểm số hoàn
toàn vô dụng làm dụng cụ đo rò rỉ mà vẫn tắt trên mọi mô hình cứng, rồi cho biết
lập luận thứ nhất và lập luận thứ ba loại nó ở bước nào.
  \item Đọc bảng~\ref{tab:ch15-tuong-quan-diem-so} và phát biểu lại đúng câu mà
bài báo nói quá, tức viết một câu về quan hệ giữa CTL, ICL và OIS mà cả tám dòng
của bảng đều ủng hộ. Sau đó cho biết nếu chỉ giữ ba dòng đầu, tức ba bộ dữ liệu
có khái niệm ít tương quan, thì câu ấy phải sửa thế nào.
  \item Điểm số can thiệp được mô tả là đặc hiệu hoàn hảo mà không nhạy. Hãy nêu
một mô hình có $\mathrm{CTL}$ và $\mathrm{ICL}$ đều bằng 0 mà hiệu năng khi can
thiệp vẫn tụt, rồi cho biết chỗ tụt ấy đến từ đâu và vì sao nó không mâu thuẫn
với việc gọi điểm số can thiệp là đặc hiệu.
  \item Mục~\ref{sec:ch15-phan-con-lai} xếp chỗ yếu mà khảo sát gọi là
\enquote{concept fidelity} vào dòng tính ổn định của
bảng~\ref{tab:ch01-bon-quan-he}. Hãy làm điều tương tự cho ba chỗ yếu còn lại
của nhánh không giám sát, và với chỗ nào không xếp vào dòng nào được thì nói rõ
bảng ấy thiếu quan hệ gì.
  \item Mục~\ref{sec:ch15-doc-duoc-dieu-khien-duoc} chỉ ra rằng cùng một công cụ
dò tên vừa gán tên vừa cung cấp cả hai điểm số. Hãy đề xuất một cách đo khả năng
điều khiển không dùng tới cái tên ấy, rồi cho biết cách của bạn đo được tính
chất nào trong hai tính chất và mất tính chất nào.
  \item Tỉ lệ 18.84 phần trăm là tỉ lệ nơ-ron đạt cả hai tính chất. Hãy cho biết
con số ấy chặn được điều gì trong ba điều sau và không chặn được hai điều kia:
mức quan trọng của một đơn vị, độ trung thực theo
định nghĩa~\ref{def:ch01-do-trung-thuc}, và độ khớp giữa tên gọi của một đơn vị
với vai trò thật của nó.
  \item Đọc mục 9, \enquote{Lessons for model design}, trong \enquote{Leakage
and Interpretability in Concept-Based Models} của Parisini và cộng
sự~\autocite{p30leakage}. Quy trình ở đó bắt đầu bằng việc huấn luyện vài đầu
phân loại trên khái niệm thật. Hãy cho biết bước ấy đo được nguyên nhân nào
trong bốn nguyên nhân của mục~\ref{sec:ch15-nguyen-nhan}, vì sao nó không tách
được hai nguyên nhân sau khỏi nhau, và cần thêm phép đo nào để tách.
\end{cauhoids}
